\documentclass[aspectratio=169]{beamer}

% Usar o estilo personalizado Newton Paiva
\usepackage{../styles/newton-paiva-style}

% Pacotes úteis
\usepackage[utf8]{inputenc}
\usepackage[portuguese]{babel}
\usepackage{graphicx}
\usepackage{hyperref}
\usepackage{enumerate}
\usepackage{amsmath}
\usepackage{amsfonts}

% Informações da apresentação
\title{Apresentação Newton Paiva}
\subtitle{Formato 16:9 com Logo no Topo}
\author{Seu Nome}
\institute{Centro Universitário Newton Paiva}
\date{\today}

\begin{document}

% Slide de título com logo no topo
\begin{frame}
    \titlepage
\end{frame}

% Sumário
\begin{frame}{Sumário}
    \tableofcontents
\end{frame}

% Seções da apresentação
\section{Formato 16:9}
\begin{frame}{Vantagens do Formato 16:9}
    \begin{columns}
        \begin{column}{0.6\textwidth}
            \begin{itemize}
                \item Mais espaço horizontal
                \item Melhor aproveitamento da tela
                \item Formato moderno e profissional
                \item Ideal para apresentações corporativas
            \end{itemize}
        \end{column}
        \begin{column}{0.4\textwidth}
            \begin{center}
                \textcolor{white}{\large\textbf{16:9}}
                \vspace{0.5cm}
                
                \textcolor{white}{vs}
                \vspace{0.5cm}
                
                \textcolor{white}{\large\textbf{4:3}}
            \end{center}
        \end{column}
    \end{columns}
\end{frame}

\section{Logo Posicionado}
\begin{frame}{Logo no Cabeçalho}
    \begin{block}{Posicionamento Estratégico}
        O logo da Newton Paiva está posicionado no topo direito de todos os slides, garantindo visibilidade constante da marca.
    \end{block}
    
    \begin{alertblock}{Identidade Visual}
        Mantém a identidade visual da instituição em destaque durante toda a apresentação.
    \end{alertblock}
    
    \begin{exampleblock}{Design Responsivo}
        O layout se adapta perfeitamente ao formato widescreen 16:9.
    \end{exampleblock}
\end{frame}

\section{Demonstração}
\begin{frame}{Conteúdo Aproveitando o Espaço}
    \begin{columns}
        \begin{column}{0.33\textwidth}
            \textbf{Coluna 1}
            \begin{itemize}
                \item Item A
                \item Item B
                \item Item C
            \end{itemize}
        \end{column}
        \begin{column}{0.33\textwidth}
            \textbf{Coluna 2}
            \begin{enumerate}
                \item Primeiro
                \item Segundo
                \item Terceiro
            \end{enumerate}
        \end{column}
        \begin{column}{0.33\textwidth}
            \textbf{Coluna 3}
            \begin{center}
                \textcolor{white}{\textbf{Destaque}}
                
                \vspace{0.5cm}
                \newtonlogo[0.8cm]
            \end{center}
        \end{column}
    \end{columns}
    
    \vspace{1cm}
    \begin{center}
        \textcolor{white}{\textbf{Formato 16:9 permite organizar melhor o conteúdo}}
    \end{center}
\end{frame}

\section{Conclusão}
\begin{frame}{Resultados Obtidos}
    \begin{itemize}
        \item ✅ Formato 16:9 implementado com sucesso
        \item ✅ Logo posicionado no topo de todos os slides
        \item ✅ Fundo personalizado mantido (\texttt{\#5c2438})
        \item ✅ Layout otimizado para apresentações modernas
        \item ✅ Identidade visual Newton Paiva preservada
    \end{itemize}
    
    \vspace{1cm}
    \begin{center}
        \textcolor{white}{\Large\textbf{Pronto para apresentações profissionais!}}
    \end{center}
\end{frame}

% Slide de agradecimentos
\begin{frame}
    \begin{center}
        \Huge\textcolor{white}{Obrigado!}
        \vspace{1cm}
        
        \Large\textcolor{white}{Perguntas?}
        
        \vspace{1cm}
        \textcolor{white}{\small Newton Paiva - Transformando o futuro}
    \end{center}
\end{frame}

\end{document}
